\cotefile{93}
\scanpage{165}
\[
 H^*_{\mathrm{dR}}(X/T)\Longleftarrow E_2^{pq}
 =H^p(S,\alpha\longmapsto\mathscr H^q(\Omega^\alpha_{S/T}
 \overset{L}{\otimes}\mathbb R f_*(\Omega^\cdot_{X/S}))),
 \tag*{(12)}
\]
où l'opérateur diff. pour le complexe $\alpha\longmapsto\ldots$ est celui déduit de la connexion de Gauss-Manin de $\mathbb R f_*(\Omega^\cdot_{X/S})$. Dans le cas favorable $\Omega^1_{S/T}$ plat, \inserted{(il suffit $\Omega^1_{S/T}$ plat, par D.~Lazard)} \deleted{\unsure{la flèche précédente}} la suite spectrale devient
\[
\begin{aligned}
 H^*_{\mathrm{DR}}(X/T)\Longleftarrow E_2^{pq}
 &=H^p(S,\alpha\longmapsto\Omega^\alpha_{S/T}\otimes\mathbb R^qf_*(\Omega^\cdot_{X/S}))\\
 &=H^p_{\mathrm{DR}}(S/T,\mathbb R^q f_*(\Omega^\cdot_{X/S})),
\end{aligned}
\tag*{(12')}
\]
où l'expression $H^p_{\mathrm{DR}}$ a un sens, grâce à la connexion de Gauss-Manin de $\mathbb R^qf_*(\Omega^\cdot_{X/S})$.

Comme de juste, \unsure{les suites spectrales|la suite spectrale} (12) \deleted{et (12')} (donc (12')) s'appelle \uline{suite spectrale de Leray} \deleted{\unsure{pour}} \unsure{de} $f$, en cohomologie de De Rham.

\deleted{\unsure{Voici une}} \uline{Application}\quad Soit $n$ un entier, et soit
\[
 \xi\in H^n_{\mathrm{DR}}(X/T)
 \tag*{(13)}
\]
telle que son image dans $E_2^{0,n}=\deleted{\unsure{H|\mathbb H^0}}\operatorname{Ker}(H^0(S,\mathbb R^n f_*(\Omega^\cdot_{X/S}))\deleted{\to H^0}\inserted{\to H^0(S,\mathscr H^n(\Omega^1_{X/S}\overset{L}{\otimes}\mathbb R f_*(\Omega^\cdot_{X/S})))})$ \unsure{par}
\[
 H^n_{\mathrm{DR}}(X/T)\longrightarrow H^0(S,\underbrace{\mathbb R^n f_*(\Omega^\cdot_{X/S})}_{\mathscr H^n_{\mathrm{DR}}(X/S)}),
 \quad\deleted{\unsure{(\to H^0(S,\mathscr H^n\mathbb R f_*(\Omega^\cdot_{X/S}))) }}
 \tag*{(14}
\]
est \uline{nulle} [\uline{NB} cette image est en tous cas tjrs \uline{horizontale}, i.e. son image dans $\Gamma(S,\mathscr H^n(\Omega^1_{X/S}\overset{L}{\otimes}\mathbb R f_*(\Omega^\cdot_{X/S})))$ est nulle). 
% archive sheet number: 164

\scanpage{166}
Alors grâce à la suite spectrale de Leray, $\unsure{\zeta|\xi}$ définit un élément dans $E_\infty^{1,n-1}\subset E_2^{1,n-1}$, donc un
\[
 \eta\in\mathbb H^1(S,\alpha\longmapsto\mathscr H^{n-1}
 (\Omega^\alpha_{S/T}\overset{L}{\otimes}\mathbb R f_*\Omega^\cdot_{X/S})).
 \tag*{(15)}
\]
Supposons en particulier les $\Omega^1_{S/T}$ plats, \deleted{de sorte} alors on obtient
\[
 \eta\in {\unsure{\mathbb H|H}}^1_{\mathrm{DR}}(S,\mathbb R^{n-1}f_*(\Omega^\cdot_{X/S}))
 \qquad(\Omega^1_{S/T}\text{ plat}).
 \tag*{(16)}
\]
Nous avons vu $[\,]$ qu'un tel objet s'identifie à une classe d'isomorphie de torseurs \inserted{$P$} sous $\mathbb R^{n-1}f_*(\Omega^\cdot_{X/S})$ muni d'un hom
\[
 P\xrightarrow{d_P}\Omega^1_{S/T}\otimes\mathbb R^{n-1}f_*(\Omega^\cdot_{X/S})
 \tag*{(17)}
\]
satisfaisant
\[
 \left\{\begin{aligned}
 d_P(x+\xi)&=d_P(x)+d^0(\xi)
 &\quad&\left\{\begin{aligned}
 x&\in P(\cdot),\quad \deleted{\unsure{\zeta}}\xi\in\mathbb R^{n-1}f_*(\Omega^\cdot_{X/S})(\cdot)\\
 d^0&\text{ déduit de conn. de G.M.}
 \end{aligned}\right.\\
 \operatorname{Im}d_P&\subset Z^1(\Omega^\cdot_{S/T}\otimes\mathbb R^{n-1}f_*(\Omega^\cdot_{X/S})).
 \end{aligned}\right.
 \tag*{(18)}
\]

[On sait aussi, grâce à l'hypothèse de platitude sur $\Omega^1_{S/T}$, que si $S$ est de car. nulle, la donnée de $d_P$ équivaut à la donnée de
\[
 P\xrightarrow{\delta_P^\infty}P_{S/T}^\infty(\mathbb R^{n-1}f_*(\Omega^\cdot_{X/S}))^+
 \mathrel{\vcenter{\hbox{\tikz{\draw (.09,.14)--(0,0)--(.22,0)--(.12,-.16);}}}}
 \tag*{(19)}
\]
satisfaisant la condition
\[
 \delta_P^\infty(x+\zeta)=\delta_P^\infty(x)+d_{S/T}^+(\zeta)
 \qquad x\in P(\cdot),\ \zeta\in\mathbb R^{n-1}f_*(\Omega^\cdot_{X/S})(\cdot)
 \tag*{(20)}
\]
\margin{$d^+_{S/T}=d^\infty_{S/T}-\delta^\infty_{S/T}$\\
$\swarrow$\\
$\delta^\infty_{S/T}$ défini par la stratification de $\mathbb R^{n-1}f_*(\Omega^\cdot_{X/S})$ déduite de sa connexion de G.M.}
\deleted{\unsure{(21)}}
% archive sheet number: 165

\scanpage{167}
\uline{Variantes}\qquad a) On peut prendre sur $\deleted{X}$ un module de coeff. qu. coh. $E$, avec connexion intégrable rel. à $T$, d'où $\Omega^\cdot_{X/T}\otimes E$ avec structure de complexe. Tout ce qui précède s'applique, (grâce au fait que les suites exactes (4) splittent localement, donc \deleted{\unsure{restent exactes en tensorisant les filtrés|suite exacte en tensorisant les filtrations}} de même pour la filtration correspondante sur $\Omega^\cdot_{S/T}\otimes E$, donc
\[
 \mathrm{Gr}_{X/S}^{p}(\Omega^\cdot_{X/T}\otimes E)
 \simeq\mathrm{Gr}_{X/S}^{p}(\Omega^\cdot_{X/T})\otimes E
\]
du point de vue additif. Pour pouvoir appliquer la formule de projection, et obtenir, en analogie avec \inserted{(3)} \deleted{\unsure{(3)}},
\[
 \mathrm{Gr}_{X/S}^{p}(\mathbb R f_*(\Omega^\cdot_{X/T}\otimes E))
 \simeq\Omega^p_{S/T}\overset{L}{\otimes}\mathbb R f_*(\Omega^\cdot_{\unsure{Y|X}/S}\otimes E)[-p],
 \tag*{(21)}
\]
il semble qu'il faille supposer $E$ \uline{plat sur} $X$ \uline{ou} $\Omega^1_{S/T}$ \uline{plat sur} $S$. Sous cette réserve, on obtient des \deleted{stratifications} connexions intégrables sur les $\mathbb R f_*(\Omega^\cdot_{X/S}\otimes E)$ rel. à $T$, et une suite spectrale qui, dans le $2^{\text{ième}}$ cas ($\Omega^1_{S/T}$ plat) prend la forme simple
\[
 H^*_{\mathrm{DR}}(X/T,E)\Longleftarrow E_2^{pq}=H^p_{\mathrm{DR}}(S/T,\mathscr H^q_{\mathrm{DR}}(X/S,E))
 \tag*{(22)}
\]

b) On peut localiser la suite spectrale obtenue sur $T$, (plus généralement, relativement à un morphisme quelconque $S\longrightarrow R$). On peut dans ce cas \deleted{\illegible{1}}\inserted{formuler} des compatibilités de la suite spectrale obtenue sur $T$, avec des connexions intégrables relativement à une base en dessous de $T$.
% archive sheet number: 166

\scanpage{168}
\begin{flushright}
\uline{Opération de Cartier}
\end{flushright}
% Exposed end of a line on the sheet beneath the cover, at the right edge.
\margin{\unsure{nt}}
% archive sheet number: 167

\scanpage{169}
\uline{L'opération de Cartier}\qquad (NB Tous les anneaux sont des $\mathbb F_p$-alg.)

\textcircled{1} Cohomologie de De Rham d'une algèbre admettant une $p$-base.

a) Soit $B/A$, alors, posant $\Lambda=A[B^p]\subset B$, on a
\[
 \left\{\begin{aligned}
 \Omega^1_{B/A}&\xrightarrow{\sim}\Omega^1_{B/\Lambda}
 &\text{d'où}\quad\Omega^*_{B/A}&\xrightarrow{\sim}\Omega^*_{B/\Lambda}\\
 \text{donc}\quad H^*_{\mathrm{DR}}(B/A)&\simeq H^*_{\mathrm{DR}}(B/\Lambda)
 &&\text{(comme $A$-algèbres)}
 \end{aligned}\right.
\]

b) Supposons que $B$ admette une $p$-base sur $A$ formée d'\uline{un} élément $x$, i.e., \deleted{\unsure{comme les $\Lambda$}} posant $x^p=a\in\Lambda$, on a un isom. de $\Lambda$-algèbres
\[
 B\simeq\Lambda[X]/(X^p-a)
\]
Alors on a
\diagram{\begin{tikzcd}[ampersand replacement=\&,column sep=.25em,row sep=1.8em]
 \text{base de }\Omega^0_{B/\Lambda}\text{ sur }\Lambda:
 \& {} \& 1,\arrow[dl] \& x,\arrow[dl,dash] \& x^2,\arrow[dl,dash,"2"'] \& \ldots
 \& x^{\unsure{p-2|p-1}}, \& x^{p-1}\arrow[dl,dash]\\
 \text{--- }\Omega^1_{B/\Lambda}\text{ --- }\Lambda:
 \& 0 \& dx, \& x\,dx, \& x^{\unsure{1|2}}dx \& \ldots \& {} \& x^{p-1}dx
\end{tikzcd}}
d'où base de \deleted{\unsure{$\Omega$}} $H^*_{\mathrm{DR}}(B/\Lambda)\simeq H^{\unsure{\alpha|*}}_{\mathrm{DR}}(B/A):\qquad\dot{1},\quad\overset{\cdot}{\overline{x^{p-1}dx}}$

c) Supposons qu'on ait une $p$-base $(x_i)_{i\in I}$ de $B$ sur $A$. Alors \deleted{\illegible{1}}
\[
 \Omega^1_{B/A}\simeq\Omega^1_{B/\Lambda}\quad\text{admet base}\quad (dx_i)_{i\in I}
\]
d'où
\[
 \Omega^*_{B/\Lambda}\simeq{\mathop{\bigotimes}\limits_I}_{\Lambda}\,\Omega^*_{B_i/\Lambda}
 \qquad\text{isom de $\Lambda$-alg différentielles}
\]
où
\[
 B_i\simeq\Lambda[x_i]\simeq\Lambda[X_i]/(X_i^p-a_i)
 \qquad (a_i=x_i^p\in\Lambda).
\]
Par suite, Künneth nous donne
\[
 H^*_{\mathrm{DR}}(B/A)\simeq H^*_{\mathrm{DR}}(B/\Lambda)
 \simeq{\mathop{\bigotimes}\limits_I}_{\Lambda}\,H^*_{\mathrm{DR}}(B_i/\Lambda)
\]
% archive sheet number: 168

\scanpage{170}
\textcircled{2} L'homomorphisme
\[
 \psi:\Omega^*_{B/A}\longrightarrow Z^*_{B/A}\qquad
 (\text{cycles de }\Omega^*_{B/A})
\]
Il y a un hom unique $\psi$, satisfaisant les conditions

a) $\psi$ est un hom d'anneaux

b) $\psi_0(x)=\underbrace{x^p}_{\unsure{F_B(x)|\psi_0(x)}}$ si $x\in B\simeq\Omega^0_{B/A}$.

c) $\psi(d_{B/A}x)=\underbrace{x^{p-1}d_{B/A}x}_{\partial x}$ si $x\in B$.

Pour vérifier ceci, il suffit de vérifier que l'on a
\[
 \partial(xy)=\psi_0(x)\partial y+\psi_0(y)\partial x
\]
i.e.
\[
 (xy)^{p-1}d(xy)=x^py^{p-1}dy+y^px^{p-1}dx
\]
ce qui est trivial.

En considérant $Z^*_{B/A}$ comme une algèbre sur $\Lambda\inserted{=A[B^p]}$, et par suite sur $\widetilde B=B\otimes_A\widetilde A$ ($\widetilde A=(A,F_A)$, $F_A:A\longrightarrow A$ donné par $x\mapsto x^p$),
\diagram{\begin{tikzcd}[ampersand replacement=\&]
 B\arrow[r]\arrow[rr,bend left=22,"F_A"] \& \widetilde B=B\otimes_A\widetilde A\arrow[r]\arrow[d] \& B\\
 A\arrow[u]\arrow[r] \& \widetilde A=(A,F_A)\arrow[ur,dash] \& {}
\end{tikzcd}}
dont l'image dans $B$ est précisément $\Lambda=A[B^p]$. Alors, utilisant l'hom canonique $B\longrightarrow\widetilde B$, $x\longmapsto x\underset A\otimes1_{\widetilde A}$, dont le composé avec $\widetilde B\longrightarrow B$ est $F_B:x\longmapsto x^p$, et considérant $\Omega^*_{B\unsure{/A}}$ comme une $B$-algèbre, l'hom. $\psi$ est \unsure{com}\illegible{1} avec l'hom $B\longrightarrow\widetilde B$ sur les scalaires, \unsure{d}\illegible{1} un \uline{hom. de $\Lambda$-algèbres}
\[
 \widetilde\psi:\Omega^*_{B/A}\otimes_B\widetilde B\simeq
 \Omega^*_{\widetilde B/\widetilde A}\longrightarrow Z^*_{\unsure{B}}
\]
\deleted{\unsure{qui a donc une base|qui est donc \illegible{1}}} Nous serons intéressé surtout par \unsure{l'hom} composé avec $Z^{\unsure{\alpha|*}}_{B/A}\longrightarrow H^\alpha_{\mathrm{DR}}(B/A)$ savoir
\[
 \boxed{\varphi:\Omega^*_{\widetilde B/\widetilde A}\longrightarrow H^*_{\mathrm{DR}}(B/A)}
\]
% Lower right of the sheet is torn away, across several line ends.
% archive sheet number: 169

\scanpage{171}
\textcircled{3} L'opération de Cartier.

\uline{Théorème 1} Supposons que $B/A$ admette une $p$-base. Alors l'hom précédent $\varphi$ \deleted{est}\inserted{induit} un \uline{isomorphisme} $\Omega^*_{\widetilde B/\widetilde A}\otimes_{\widetilde B}\Lambda\longrightarrow H^{\unsure{\alpha|*}}_{\mathrm{DR}}(B/A)$ [NB $\Lambda\simeq\widetilde B/$\unsure{idéal|ident}] donc est un isom si $\widetilde B\xrightarrow{\sim}\Lambda$ i.e. $\widetilde B\longrightarrow B$ est injectif.

En effet, si $(x_i)_{i\in I}$ est une $p$-base, utilisant les notations du n$^\circ$ 1, on trouve \deleted{\unsure{un}}\inserted{un isom, de $B$-algèbres}
\diagram{\begin{tikzcd}[ampersand replacement=\&,column sep=1em,row sep=.75em]
 B\arrow[r]\arrow[d] \& \widetilde B\arrow[r]\arrow[d,dash] \& B\arrow[d]\\
 B_i\arrow[r,dash]\arrow[d] \& \widetilde B_i\arrow[r,dash]\arrow[d,dash] \& B_i\arrow[d]\\
 \Lambda\arrow[r,dash]\arrow[d,dash] \& \widetilde\Lambda\arrow[r,dash]\arrow[d,dash] \& \Lambda\\
 A\arrow[r,dash] \& \widetilde A \& {}
\end{tikzcd}}
\[
 \Omega^*_{B/A}\simeq\bigotimes_{i\in I}\Omega^*_{B_i/\Lambda}\otimes_{B_i}B
 \quad\text{d'où}\quad
 \Omega^\alpha_{\widetilde B/\widetilde A}\simeq\bigotimes_{i\in I}
 (\overset{\deleted{\illegible{1}}}{\Omega^*_{\widetilde B_i/\widetilde\Lambda}}\otimes_{\unsure{B_i|\widetilde B_i}}\widetilde B),
\]
\margin{isom de $B$-alg.}
et en vertu du calcul du n$^\circ$ 1 pour $H^\alpha_{\mathrm{DR}}(B/A)$ :
\[
 H^*_{\mathrm{DR}}(B/A)\simeq\bigotimes_{i\in I}H^\alpha_{\mathrm{DR}}(B_i/\Lambda)
\]
L'hom $\varphi$ envisagé s'identifie alors au produit tensoriel des homomorphismes analogues \deleted{\unsure{envisagés|correspondants}} de
\[
 \begin{aligned}
 \varphi_i:\Omega^*_{\widetilde B_i/\widetilde\Lambda}\otimes_{B_i}\widetilde B_i
 &\longrightarrow H^*_{\mathrm{DR}}(B_i/\Lambda)\\
 \text{base }\ \widetilde 1,\ d\widetilde x_i&\qquad
 \text{base }\ \dot{1},\ \overset{\cdot}{\overline{x_i^{p-1}dx_i}}
 \end{aligned}
\]
qui envoient base sur base, donc sont des isom. Cela prouve le théorème.

\uline{Corollaire} 1 \deleted{\unsure{Supposons de plus|Supplémentaire}} On a
\[
 Z^*(B/A)\simeq B^*(B/A)\oplus[\Omega^*_{\widetilde B/\widetilde A}\otimes_{\widetilde B}\Lambda]
\]

\uline{Définition}\quad L'isom inverse de $\varphi$ est noté
\[
 \boxed{\widetilde C:H^*_{\mathrm{DR}}(B/A)\xrightarrow{\sim}
 \Omega^*_{\widetilde B/\widetilde A}\otimes_{\widetilde B}\Lambda}
 \boxed{(=\Omega^\alpha_{B/A}\otimes_B\Lambda)}
\]
On dénote aussi souvent par $C$ l'hom.
\[
 \widetilde C:Z^*(B/A)\longrightarrow\Omega^*_{\widetilde B/\widetilde A}
\]
donnant le précédent par passage au quotient, [donc ce dernier est \uline{surjectif}, \uline{et son noyau est formé des différentielles exactes}].
% archive sheet number: 170

\scanpage{172}
\uline{Corollaire 2}\quad Il y a une application unique
\[
 \widetilde C:Z^*(B/A)\mathrel{\text{---}}\Omega^*_{\widetilde B/\widetilde A}
\]
satisfaisant aux conditions suivantes

a) $\widetilde C(d\omega)=0$ pour \unsure{tte|tout} $\omega\in Z^*(B/A)$

b) $\widetilde C$ est un hom. de \deleted{\unsure{anneaux}}\inserted{\uline{$\Lambda$-algèbres}}
\[
 \left\{\begin{aligned}
 \widetilde C(\omega+\omega')&=C(\omega)+C(\omega')\\
 \widetilde C(\omega\omega')&=C(\omega)\,C(\omega')\\
 \widetilde C(\lambda\omega)&=\lambda C(\omega)\quad\text{si }\lambda\in\Lambda.\\
 C(1)&=1
 \end{aligned}\right.
\]

c) $\widetilde C(\deleted{\unsure{\xi}}x^{p-1}d_{B/A}x)=d_{\widetilde B/\widetilde A}x\otimes_{\widetilde B}1_\Lambda$

L'existence résulte de la définition. L'unicité résulte du corollaire 1.

\uline{Corollaire 3}\quad Si $x\in B^*$, on a
\[
 \text{d) }\ \widetilde C\left(\frac{dx}{x}\right)=\frac{dx}{x}\otimes_B1_\Lambda
 =\frac{d\widetilde x}{\widetilde x}\otimes_{\widetilde B}1_\Lambda
\]
En effet, on a $\frac{dx}{x}=x^{-p}(x^{p-1}dx)$, et comme $x^{-p}\in\Lambda$, on trouve
\[
 \begin{aligned}
 \widetilde C\left(\frac{dx}{x}\right)
 &=x^{-p}\unsure{\overline C|\widetilde C}(x^{p-1}dx)=\deleted{\unsure{x\,d\widetilde x}}F_A(x^{-1})C(x^{p-1}dx)\\
 &=(\widetilde x^{-1}d\widetilde x)\otimes_{\widetilde B}\Lambda
 \end{aligned}
\]

\uline{Cas particuliers}

a) Supposons que $\operatorname{Im}(A\mathrel{\text{--}} B)\deleted{\unsure{\supset}}\subset B^p$, \deleted{\unsure{i.e. $A=\operatorname{Im}(A\to B)$, Alors|i.e. $A\to\operatorname{Im}(A\to B)$. Alors}} p.ex $A$ parfait. Alors $\Omega^*_{B/A}\simeq\Omega^\alpha_{B/\mathbb F_p}\overset{\text{défn}}{=}\Omega^\alpha_B$, et on peut pour étudier \inserted{$H^*_{\mathrm{DR}}(B/A)$} supposer que $A=\mathbb F_p$. Mais, dans ce dernier cas, on aura $A\xrightarrow{\sim}\widetilde A$, donc $B\simeq\widetilde B$, et si $B\xrightarrow{F_B}B$ est injectif i.e. $\widetilde B\xrightarrow{\sim}\Lambda$ i.e. $B$ \uline{réduit}, alors le th. 1 revient à $\Omega^*_{B\deleted{\unsure{/A}}}\xrightarrow[\sim]{\varphi}H^*_{\mathrm{DR}}(B)$. Donc l'opération de Cartier peut être considérée alors comme
\[
 \overset{\deleted{\sim}}{C}:\deleted{\unsure{B}}Z^*(B)\longrightarrow\Omega^*_B
\]
On fera attention que la \inserted{dernière} formule s'écrit ici
\[
 \text{b')}\qquad \overset{\deleted{\sim}}{C}(f^p\omega)=fC(\omega)
\]
et c) et d) deviennent
% archive sheet number: 171

\scanpage{173}
\[
 \left\{\begin{aligned}
 C(x^{p-1}dx)&=dx\\
 C(dx/x)&=dx/x
 \end{aligned}\right.
\]
C'est le cas envisagé dans la note initiale de Cartier.

\textcircled{4} Fonctorialités, Globalisation.

Soit
\diagram{\begin{tikzcd}[ampersand replacement=\&,row sep=.65em]
 X\arrow[d]\\ S
\end{tikzcd}}
en car $p$, d'où
\diagram{\begin{tikzcd}[ampersand replacement=\&,column sep=2em]
 X\arrow[d,dash] \& X^{(p/S)}\arrow[l]\arrow[d,dash] \& X\arrow[l,"{f=f_{X/S}}"']\arrow[ll,bend right=25,"F_X"']\arrow[dl,dash]\\
 S \& S\arrow[l,"F_S"]\arrow[ul,phantom,"\text{\uline{cart.}}",very near start] \& {}
\end{tikzcd}}
, avec $f=f_{X/S}$ radiciel \deleted{\unsure{entier}} surjectif.

\deleted{Supposons que $f_*(\underline O_X)$ ait \unsure{une} localement une $p$-base sur $\operatorname{Im}(\underline O_{X^{(p/S)}}\longrightarrow f_*(\underline O_X))$, alors on obtient des isom. can.}

On définit alors un hom canonique de faisceaux d'Algèbres
\[
 \psi:\Omega^*_{X^{p/S}/S}\longrightarrow f_*(Z^*(\Omega^*_{X/S}))
\]
d'où un hom composé
\[
 \psi:\Omega^*_{X^{p/S}/S}\longrightarrow f_*(\mathscr H^*(\Omega^*_{X/S}))
 \simeq\mathscr H^\alpha(f_*(\deleted{\unsure{Z}}\!(\Omega^*_{X/S}))),
\]
Si $\underline\Lambda=\operatorname{Im}(\underline O_{X^{p/S}}\longrightarrow f_*\underline O_X))$ [qui est donc de la forme $\underline O_{X^{p/S}}/J$, $J$ un nilidéal, donc $J=0$ si $X^{p/S}$ est réduit, p.ex. $S$ réduit et $X$ \deleted{\unsure{lisse}}\inserted{loc} de présentation finie et \unsure{ouvert|u. ouvert} sur $S$], on aura donc
\[
 \widetilde\psi:\Omega^\alpha_{X^{p/S}/S}\otimes_{\underline O_{X^{p/S}}}\underline\Lambda
 \longrightarrow\deleted{f_*}(\mathscr H^\alpha(f_*(\Omega^*_{X/S})))
\]
Si $f_*(\underline O_X)$ admet loc une $p$-base sur $\underline O_{X^{p/S}}$ (ou ce qui revient au même, sur $\Lambda$), alors l'hom précédent est un isom, et on peut définir, \deleted{\unsure{la ca}} l'isom inverse
\[
 \widetilde C:\mathscr H^\alpha(f_*(\Omega^\cdot_{X/S}))
 \longrightarrow\Omega^*_{X^{p/S}/S}\otimes_{\underline O_{X^{p/S}}}\Lambda
\]
\margin{définir $\widetilde C$ chaque fois que $\psi$ est un isom (sans qu'il y ait néc. des $p$-bases\ldots)}
d'où un hom
\[
 \widetilde C:Z^*(f_*(\Omega^*_{X/S}))\longrightarrow
 \Omega^*_{\deleted{\unsure{X/S}}X^{p/S}/S}\otimes_{\underline O_{X^{p/S}}}\Lambda
\]
% archive sheet number: 172

\scanpage{174}
La formation de $\psi$, et de $C$ sous réserve d'existence, satisfait à une propriété de fonctorialité évidente pour un carré commutatif
\diagram{\begin{tikzcd}[ampersand replacement=\&]
 X\arrow[d] \& X'\arrow[l]\arrow[d]\\
 S \& S'\arrow[l]
\end{tikzcd}}
donnant naissance à un diagramme comm.
\diagram{\begin{tikzcd}[ampersand replacement=\&,row sep=1.2em,column sep=1.3em]
 {} \& X\arrow[dl]\arrow[dd] \& {} \& X'\arrow[ll]\arrow[dl]\arrow[dd]\\
 X^{p/S}\arrow[dr] \& {} \& {X'}^{p/S'}\arrow[ll]\arrow[dr] \& {}\\
 {} \& S \& {} \& S'\arrow[ll]
\end{tikzcd}}
\noindent ---\,---\,---\,---

\textcircled{5} \deleted{Relations avec les} \uline{Connexions sur les Modules inversibles, et différentielles logarithmiques}

Soit $\underline L$ module inversible sur $X$ \inserted{muni d'une connexion $\Delta_0$ rel. $S$.} \deleted{\unsure{inversible}} Alors les connexions \inserted{$\Delta$} sur $\underline L$, \deleted{qui} (rel. à $S$), \deleted{s'identifient}\inserted{qui correspondent} aux hom linéaires $u:\underline L\longrightarrow\underline L\otimes\Omega^1_{X/S}$ \inserted{par $u\longmapsto\Delta=\Delta_0+u$} s'identifient aux 1-formes différentielles \deleted{\unsure{les}} $\omega$ sur $X/S$.

\uline{Proposition}\quad \inserted{Si $\Delta_0$ est à courbure nulle,} Pour que la connexion sur $\underline L$ définie par $\omega$ soit \deleted{\unsure{sans torsion}}\inserted{à courbure nulle}, il f. et s. que l'on ait $d\omega=0$.

\deleted{Pour que \unsure{de plus|elle} soit} Supposons que $X/X^{p/S}$ ait loc une $p$-base, et que \deleted{\unsure{désormais}}\inserted{$\Delta,\Delta_0$ soient à courbure nulle, alors $\Delta_0$ comp. aux puiss. $p^{\text{ièmes}}$} Pour que la connexion envisagée commute aux puissances $p^{\text{ièmes}}$ [i.e. $\Theta_\xi^p(x)=\Theta_{\xi^p}(x)$ pour tt champ de vecteurs $\xi$ sur un ouvert, et \deleted{\unsure{section}} $\unsure{a|x}$ \deleted{\unsure{du|de}} $\underline L$ sur celui-ci)
% archive sheet number: 173

\scanpage{175}
il faut et il suffit que l'on ait $\widetilde C\omega=\widetilde\omega\otimes_{\underline O_{X^{p/S}}}1_{\underline\Lambda}$.

La première assertion est \unsure{standart|standard}, et ne dépend pas de l'hyp. de car. $p>0$. [Si $\underline L$ n'est pas de rang 1, mais seulement loc. libre, alors $\omega$ est une forme diff. à valeurs dans $\underline{\operatorname{End}}(E)$, et la relation à écrire est $d\omega+[\omega,\omega]=0$]

% slanted
\margin{\uline{maladroit de choisir des bases} Dire : la section de $(\Omega^1_{X/S})^{\otimes p}$ ``défaut de compatibilité aux puiss. $p^{\text{ièmes}}$'' de la connexion envisagée $\Delta$ [ou $p$-courbure] est égale à [celle de $\Delta_0$, supposée nulle, plus] $\omega^{\otimes p}-f_{X/S}^*\widetilde C\omega=f_{X/S}^*(\widetilde\omega-\widetilde C\omega)$ donc est nulle ssi on a $\widetilde\omega=\widetilde C\omega$ \ldots.}

Pour la deuxième, la condition étant locale, on est ramené au cas où tout est affine. On a donc, si on choisit une $p$-base $(x_i)$, d'où des $D_i=\frac{\partial}{\partial x_i}$,
\deleted{$\unsure{D_\xi,\rho_\Delta}(\unsure{D})\quad\rho(D_i)=\rho_0(D_i)+\unsure{u}$}
\deleted{d'où} Ecrivons la condition que $\Delta$, \deleted{\unsure{associé à}}\inserted{qui donne} $\xi\longrightarrow\Theta_\xi$ \deleted{$\Theta(D_i)^p$} est compatible aux puissances $p$-ièmes. On trouve \deleted{$\unsure{\Theta_{D_i}^p}$} les conditions nécessaires
\[
 \Theta_{D_i}^{p}=0\qquad i\in \unsure{I|J}
\]
car $D_i^p=0$. Je dis qu'elle est suffisante. \deleted{\unsure{En}}\inserted{Cela} se voit aisément, en prouvant que l'algèbre des \deleted{\unsure{effet, tte dérivation est de la forme $\sum a_iD_i$,}} op. différentiels engendrés par les $D_i=\frac{\partial}{\partial x_i}$ et les mult. \deleted{les $a_i\in B$.} scalaires est engendrée par \inserted{les générateurs $D_i,x_i$, avec} les relations
\margin{Indépendant de l'hyp. que $L$ de rang 1}
\[
 \left\{\begin{aligned}
 [D_i,D_j]&=0\\
 [x_i,x_j]&=0\\
 [D_i,x_j]&=\delta_{ij}\\
 x_i^p&=a_i\\
 D_i^p&=0
 \end{aligned}\right.
\]
% slanted
\margin{\uline{mieux} : définir l'obstruction à la compatibilité aux puiss. $p$.ièmes comme une section de $(\Omega^1_{X/S})^{\otimes p}\otimes\underline{\operatorname{End}}(E)$ (pour une connexion \uline{à courbure nulle} donnée)}

Il faut donc, compte tenu de $\overset{\Delta_i}{\Theta_{D_i}}=\overset{\Delta_i^0}{\Theta^0_{D_i}}+\omega(D_i).\mathrm{id}$ écrivant $\omega=\sum c_i\,dx_i$, où $c_i=\omega(D_i)$, exprimer la condition
\[
 (\Delta_i^0+c_i\cdot)^p=0
\]
Mais la formule de Jacobson donne
% archive sheet number: 174

\scanpage{176}
(cf p.ex. Sém. Chevalley 58/59, exp. 6 p. 06)
\[
 (\Delta_i^0+c_i)^p=\underset{\substack{\Vert\\0\ \text{par hyp}}}{\Delta_i^{0\,p}}+c_i^p+D_i^{p-1}c_i
\]
donc la condition envisagée s'écrit
\[
 c_i^p+D_i^{p-1}c_i=0\qquad i\in I
 \tag*{(*)}
\]
Or écrivons que \inserted{$d\omega=0$,} $C\omega=\widetilde\omega$, \deleted{on trouve} $(\omega=\sum c_i\,dx_i)$ on trouve pour $d\omega=0$ la condition
\[
 D_i c_j-D_j c_i=0\qquad i,j\in I
 \tag*{(**)}
\]
et on peut alors écrire
\[
 \widetilde C\omega=
 \deleted{\sum_i[(c_i^p\unsure{+|-}D_i^{p-1}c_i)\otimes_B1_\Lambda]\,d\widetilde x_i}\ 
 \sum_i-\underbrace{D_i^{p-1}c_i}_{\substack{\text{est dans }\Lambda\\
 \text{car égal à}\\\langle\widetilde D_i,C\omega\rangle1_\Lambda}}
 (d\widetilde x_i\otimes1_\Lambda
 \tag*{(***)}
\]
\deleted{ce qui signifie aussi que $\sum(c_i^p+D_i^{p-1}c_i)x_i^{p-1}dx_i\overset{\text{cohomologue}}{\sim}\sum c_i\,dx_i,$ i.e. que $\sum c_i^p$}
[grâce à la formule de \unsure{transposition|transformation} de loc cit. p. 05 (prop. 3)].
% slanted
\margin{inclure dans le formulaire du \S\,3 ou \S\,4}

alors qu'on a $\widetilde\omega=\sum c_i^p(d\widetilde x_i\otimes1_\Lambda)$

et on trouve donc,
\[
 -\widetilde C\omega+\overline\omega=\sum(c_i^p+D_i^{p-1}c_i)(d\widetilde x_i\otimes1_\Lambda)
 \tag*{$*$}
\]
donc c'est nul ssi (*) l'est, cqfd.

Nous utilisons maintenant le fait bien connu qu'une connexion à courbure nulle \deleted{\unsure{pour|dont}} et compatible à la puissance $p$-ième, sur un Module quasi-cohérent, équivaut à une \uline{donnée de descente} de celui-ci \deleted{\unsure{de}} $X$ à $(X^{p/S},\underline\Lambda)$. \deleted{On trouve que}
% archive sheet number: 175

\scanpage{177}
\deleted{\illegible{1} loc. sur $X^{(p/S)},\ \underline O_X$} Prenons $L=\underline O_X$, avec $\Delta_0$ la connexion habituelle. Alors $\omega$ satisfaisant $d\omega=0$, $\widetilde C\omega=\widetilde\omega$ définit donc une donnée de descente sur $L$, qui \deleted{qui} provient de $L_0$ sur $X_0^{p/S}=(X^{p/S},\underline\Lambda)$. Localement, on peut alors \deleted{\unsure{splitter cette connexion,}} trouver une base de $L_0$, i.e. loclt on peut engendrer $L$ par une section horizontale $f$, ce qui signifie que loc. on a $\omega=\frac{df}{f}$. On trouve ainsi le

\uline{Théorème}\quad Pour que \inserted{la 1-forme} $\omega$ soit loc. une dérivée logarithmique, $\omega=df/f$ ($f\in\Gamma(O_X^*)$) il f. et s. que l'on ait $d\omega=0$, \inserted{$C\omega=\widetilde\omega$.} En d'autres termes, on a un hom $f\longrightarrow df/f$ :
\[
 O_X^*\longrightarrow\underline\Omega^1_{X/S}
\]
dont l'image est le sous-faisceau de $\underline\Omega^1_{X/S}$ formé des $\omega$ telles que $d\omega=0$, $\widetilde C\omega=\widetilde\omega$. [NB le noyau est $\deleted{\unsure{\underline O}}\Lambda^*$, donc si $S$ est parfait, c'est $\underline O_X^{*\,p}$]

\textcircled{6.} \uline{Cas des schémas relatifs lisses. Relations avec la trace des formes diff. de degré maximum.}

Supposons $X$ lisse sur $S$, donc $X^{p/S}$ lisse sur $S$. Alors $X\xrightarrow{f_{X/S}}X^{p/S}$ est un morphisme fini localmt libre (en plus d'être radiciel et surjectif), à
% archive sheet number: 176

\scanpage{178}
fortiori $\underline O_{X^{p/S}}\xrightarrow{\sim}\Lambda$ \deleted{\unsure{Donc On sait bi|Donc On sait bien}} [Il suffit de vérifier fibre par fibre\ldots\ \inserted{or $X_s\longrightarrow X_s$ est fini, $X_s$ (\unsure{abst|étant}) régulier, \unsure{$=$ dim pour source|$=$ dim que source|à dimension plus source}.}] On sait bien que $X/S$ a localmt une $p$-base i.e. qu'il en est ainsi de $X$ sur $X^{p/S}$ : \unsure{en effet|enf} fait, grâce au fait que $X$ est loc de t.f. sur $S$, cette condition équivaut à celle que $\Omega^1_{X/S}$ soit localement libre, ce qui est bien le cas !

On est donc dans un cas où l'opération de Cartier est
\[
 \widetilde C:f_{X/S*}^{\deleted{\unsure{*}}}(Z\Omega^*_{X/S})\longrightarrow\Omega^*_{X^{p/S}}
\]
Si la dim relative est $n$, on trouve en particulier
\[
 \widetilde C:f_{X/S*}(\underbrace{\underline\Omega^n_{X/S}}_{
 \substack{\Vert\\\omega_{X/S}\\\Vert\\\text{module}\\\text{dualisant}\\\text{relatif}}})
 \longrightarrow\underset{\substack{\Vert\\\omega_{X^{p/S}}}}{\underline\Omega^n_{X^{p/S}}}
\]

\uline{Théorème} 3\quad L'homomorphisme précédent est l'homomorphisme trace (ou de Gysin) $f_{X/S*}$ associé au morphisme $f_{X/S}:X\longrightarrow X^{p/S}$ de schémas lisses sur $S$.

Soit alors $f_1,\ldots,f_n$ une suite de sections de $\underline O_X$ qui définissent un sous-schéma fermé \inserted{$Z$} de $X$ \uline{fini} sur $S$ \inserted{(donc loc lib. sur $S$)}. $\widetilde f_1\ldots\widetilde f_n$ définissent donc $\widetilde Z$ itou, et on peut former, pour une section $\overline\omega$ de $\Omega^n_{X^{p/S}}$, $\deleted{\operatorname{Tr}}\operatorname{res}\left(\frac{\overline\omega}{\widetilde f_1\ldots\widetilde f_n}\right)$. On peut aussi regarder les images $f_1^p,\ldots,f_n^p$ dans $X$, qui définissent un sous-schéma $Z'$ de $X$ fini sur $S$, et pour les $n$-formes $\omega$ de $X/S$, on peut former $\operatorname{res}\left(\frac{\omega}{f_1^p\ldots f_n^p}\right)$. On trouve alors, comme façon essentielle
% archive sheet number: 177

